\begin{textsample}{POS Dim 3 – al – Score 15.00 – al/al\_inf\_12.txt}  \label{ex:f3_pos_173}
2.
A EDUCAÇÃO \textbf{INFANTIL} E UMA PROPOSTA DE CURRÍCULO AMPLIADO Para a construção de uma proposta curricular ampliada, que se efetive na prática, faz -se necessário que as concepções em torno dos elementos articuladores estejam bem definidos e alinhados entre todos os que o desenvolvem, especialmente em torno das especificidades que compõem a Educação \textbf{Infantil}.
Para se pensar numa proposta de currículo é necessário se definir de qual currículo estamos falando, para quem é construído, quais objetivos e como pensar a educação.
De acordo com as Diretrizes Curriculares Nacionais para a Educação \textbf{Infantil} ( DCNEI, 2009 ), defendido na BNCC, currículo é um conjunto de práticas que buscam articular as experiências e os saberes das \textbf{crianças} com os conhecimentos que fazem parte do patrimônio cultural, artístico, ambiental, científico e tecnológico, de modo a promover o desenvolvimento integral de \textbf{crianças} de 0 a 5 anos de idade.
Desse modo, este capítulo se ocupará de definir algumas principais concepções os sujeitos e os elementos que se inter-relacionam no contexto prático educativo institucional. \textbf{Brincar} é uma atividade essencialmente humana e a \textbf{brincadeira} é um espaço educativo da maior importância na infância porque alimenta o espírito imaginativo, exploratório e inventivo da \textbf{criança}.
Por meio da \textbf{brincadeira} ela aprende a expressar seus sentimentos, desenvolver a imaginação criativa, compartilhar, tomar decisões, superar problemas, observar e compreender o mundo à sua \textbf{volta}.
Além disso, contribui para o desenvolvimento da socialização, propicia o aparecimento de atitudes cooperativas e a possibilidade de experiências grupais.
Ao observar atentamente a \textbf{criança} \textbf{brincando}, o professor pode descobrir como ela vê o mundo e como ela gostaria que ele fosse, quais são suas preocupações e quais são os problemas que ela coloca para si mesma. \textbf{Brincando} ela observa, experimenta, investiga, reflete, desconstrói e reconstrói o mundo à sua maneira.
Por tudo isso, as situações de \textbf{brincadeira} contribuem significativamente para o desenvolvimento e a aprendizagem da \textbf{criança} e também favorecem que ela elabore suas emoções, conheça melhor o próprio corpo, lide com a sexualidade, ressignifique suas relações interpessoais e aprenda a se conhecer melhor.
Interação e \textbf{brincadeira} A \textbf{criança} é um ser social, que se desenvolve pela interação com outros seres humanos, portanto, ampliar possibilidades de interação, em situações de diferentes naturezas, contribui positivamente para a formação da personalidade, tornando -se essencial nesse processo de desenvolvimento \textbf{infantil}.
Por isso, a importância de, na instituição de educação \textbf{infantil}, as \textbf{crianças} experimentarem diferentes possibilidades : não apenas situações de interação com os pares da mesma faixa etária e com os profissionais \textbf{adultos} com os quais convivem regularmente, mas também situações intencionalmente planejadas ( e acompanhadas pelos professores ) de interação com \textbf{crianças} de idades diferentes, que oportunizam o intercâmbio com parceiros mais e menos experientes.
A \textbf{brincadeira} deve ocupar lugar de destaque na \textbf{rotina} das \textbf{crianças}, sendo esta uma atividade própria desta etapa da vida e exerce enorme influência sobre seu desenvolvimento, pois permite que as próprias \textbf{crianças} se organizem, privilegiando sua autonomia, socialização e sociabilidade, à medida que resulta de múltiplas relações com parceiros distintos, \textbf{livremente} escolhidos.
Quando \textbf{brincam} de faz de conta, as \textbf{crianças} representam papéis e expressam situações agradáveis ou desagradáveis e tentam, assim, dizer ao \textbf{adulto} o que estão pensando.
É nesse processo de interação que elas buscam resolver os problemas e vão desenvolvendo habilidades que lhes permite crescer e aprender a conviver com seus pares.
Assim, a proposta é trabalhar com \textbf{crianças} de várias turmas e idades compondo um único coletivo.
Conforme a concepção de \textbf{brincadeira}, defendida pelas Orientações Curriculares para a educação \textbf{infantil} da rede municipal de educação do município de Maceió ( SME, 2015, p.120-122 ), as instituições de educação \textbf{infantil} devem : - Compreender que o \textbf{brincar} contribui para a construção da identidade da \textbf{criança}, para o exercício de suas capacidades e habilidades, promove a aprendizagem e o desenvolvimento como um todo ; - Reconhecer que o contexto social da \textbf{brincadeira} é crucial para o desenvolvimento global da \textbf{criança}, que a interação entre pares é uma oportunidade importante para a aprendizagem social e que a \textbf{brincadeira} é um espaço privilegiado para a expressão das culturas \textbf{infantis} ; - Oferecer amplos espaços ( internos e externos ) para a \textbf{brincadeira} e promovê -la, diariamente, com a oferta de amplo acervo de materiais, \textbf{brinquedos} e arranjos \textbf{mobiliários} ( vide item 2.2.5 ) ; - Reconhecer o importante papel dos professores de \textbf{apoiar} a \textbf{brincadeira} da \textbf{criança} e tornar as \textbf{crianças} brincantes cada vez melhores.
Para tanto, é preciso organizar diferentes ‘ cantos’/propostas, nos quais as \textbf{crianças} se inserem conforme sua própria escolha : jogos, artes, \textbf{brincadeira} de faz de conta, etc.
Neste contexto, ao professor, cabe o papel de parceiro mais experiente, que observa, \textbf{acolhe}, propõe, intervém, acrescenta, recoloca as coisas em outras bases.
Além de organizar vários espaços com materiais diversificados - espaço de artes \textbf{plásticas}, de jogos de mesa e construção, de circuito, de biblioteca, de casinha – refletir sobre sua prática a partir da observação e análise que faz das suas atitudes e das atitudes das \textbf{crianças} durante o trabalho.
Esse tipo de prática, que poderia ser chamada de “ espaços interativos ”, favorece a constituição de procedimentos, valores e atitudes fundamentais à vida em coletividade, como a colaboração, a \textbf{ajuda} mútua, a cooperação, a solidariedade e o respeito pelo outro, por suas diferenças, além de ampliar as possibilidades de conhecimento das \textbf{crianças} sobre si mesmas e sobre os outros.
Nessa perspectiva, as culturas \textbf{infantis} vão sendo construidas baseada nos espaços de encontros e diálogos que propõe as \textbf{crianças} um modelo de \textbf{escuta} e de respeito, ao tempo em que estimula essa \textbf{criança} a buscar respostas as suas perguntas na interação com seus pares, respeitando os tempos e ritmos de cada um.

% matched lemmas: acolher, adulto, ajuda, apoiar, brincadeira, brincar, brinquedo, criança, escuta, infantil, livremente, mobiliário, plástico, rotina, volta
\end{textsample}
